\documentclass[12pt]{article}
\usepackage[utf8]{inputenc}
\usepackage[T1]{fontenc}
\usepackage{hyperref}
\title{Santhali Language Resources for K-12 Curriculum Development}
\author{Compiled by Claude Sonnet 5.5 (claude-sonnet-5-5) for Andr\'as Kornai}
\date{}
\begin{document}
\maketitle

\section{Introduction}
Santhali (ISO 639-3 sat) is an Austroasiatic language of India, Bangladesh and Nepal. It has about 7.2 million speakers and is written mainly in the Ol Chiki script. Roman, Devanagari, Bengali and Odia scripts are also in use. The script question is contested: in Assam, a Santali body has objected to Ol Chiki textbooks and asked for the Roman script. The searches below found primers, press reports, a digitized dictionary and some audio. They found no official full-grade textbook series. Several items were seen only as search snippets, and their content could not be checked.

\section{Primary Educational Resources}
\begin{itemize}
\item CIIL Santali--Odia Primer, a basal reader for the Ol Chiki alphabet and numerals, prepared by CIIL Mysuru: \url{https://ciil.org/primers/Santali-Odia_Primer.pdf}
\item CIIL/NCERT Santali--Hindi Primer, first edition May 2025. Only the front matter was visible in search, so the lesson content is unconfirmed: \url{https://www.ciil.gov.in/primers/Santali-Hindi_Primer.pdf}
\item Daily Star report on Santal children in Bangladesh not yet receiving mother-tongue textbooks: \url{https://www.thedailystar.net/backpage/santal-children-yet-get-textbooks-mother-tongue-1342030}
\end{itemize}

\section{Linguistic and Reference Materials}
\begin{itemize}
\item Digitized text of P.O. Bodding's Santal dictionary with test full-text search (page last updated 2004; the link may be dated): \url{https://www.aa.tufs.ac.jp/~mmine/GICAS/SantaliDatj-e.htm}
\item Santali language resources list (TIFR, dated 2002): \url{https://theory.tifr.res.in/bombay/history/people/language/santhali.html}
\item Omniglot page on Ol Chiki: \url{https://omniglot.com/writing/olchiki.htm}
\item Unicode proposal comment on Ol Chiki: \url{https://www.unicode.org/L2/L2002/02456-ol-chiki-comment.pdf}
\item Wikipedia on the Ol Chiki Unicode block: \url{https://en.wikipedia.org/wiki/Ol_Chiki_(Unicode_block)}
\end{itemize}

\section{Religious and Community Resources}
No Santali Bible text was confirmed in these searches. Global Recordings Network has evangelistic audio (see below). Commentary on the script dispute and school implementation:
\begin{itemize}
\item \url{https://countercurrents.org/2026/09/the-script-war-nobody-is-talking-about/}
\item \url{https://raiot.in/on-learning-and-language-a-photoessay/}
\end{itemize}

\section{Digital and Audio Resources}
\begin{itemize}
\item Global Recordings Network, Santali: \url{https://globalrecordings.net/en/language/618}
\item Global Recordings Network, Words of Life 3, Santali: \url{https://www.globalrecordings.net/en/program/63116}
\item Global Recordings Network, Santhali (Bangladesh): \url{https://www.globalrecordings.net/en/language/5030}
\item Wikimedia CIS-A2K work plan for Santali Wikipedia: \url{https://meta.wikimedia.org/wiki/CIS-A2K/Work_plan_July_2015_-_June_2016/Santali_Wikipedia}
\item IISc SPIRE Lab job description mentioning Santali speech data collection (not a resource in itself): \url{https://spire.ee.iisc.ac.in/assets/PDFs/Openings/JD-NLTM_Santali.pdf}
\end{itemize}

\section{Curriculum Development Context}
Santali is reported to be used in elementary education in Jharkhand, Odisha and West Bengal, but implementation is uneven. Ol Chiki primary textbooks were printed in West Bengal in the 1980s with mixed results. Assam introduced Ol Chiki Santali textbooks in 2025 and met objections about the script. Planners should settle the script question with the community. They should also check the official SCERT and board sources, since the practice-exam sites that turned up are third-party aids and not syllabi.
\end{document}